% !TeX program = xelatex
% UTF-8 transcription; source wording, abbreviations and numerical values retained.
% images/pic01.jpg = 20261012_194407856.jpg (page 1)
% images/pic02.jpg = 20261012_194415250.jpg (page 2)
% images/pic03.jpg = 20261012_194421436.jpg (page 3)
% images/pic04.jpg = 20261012_194427670.jpg (page 4)
% images/pic05.jpg = 20261012_194433496.jpg (page 5)
% images/pic06.jpg = 20261012_194439174.jpg (page 6)
% images/pic07.jpg = 20261012_194444740.jpg (page 7)
% images/pic08.jpg = 20261012_194450982.jpg (page 8)
% images/pic09.jpg = 20261012_194456192.jpg (page 9)
% images/pic10.jpg = 20261012_194500674.jpg (page 10)
% Standalone XeLaTeX document, or input into an existing XePersian document.
\ifdefined\settextfont
  \def\PEDVTStandalone{0}
\else
  \def\PEDVTStandalone{1}
  \documentclass[12pt,a4paper]{article}
  \usepackage[margin=2cm]{geometry}
  \usepackage{amsmath,amssymb}
  \usepackage{xepersian}
  \IfFontExistsTF{Amiri}{\settextfont{Amiri}}{\settextfont{FreeSerif}}
  \setlatintextfont{FreeSerif}
\fi
\providecommand{\en}[1]{\lr{#1}}
\newcommand{\PEDVTpage}[1]{%
  \clearpage
  \section*{موضوع: براونوالد ـ ۸۷ ـ \en{PE, DVT}}
  {\footnotesize تاریخ: \en{/ /}\quad صفحهٔ #1\par}}
\newcommand{\PEDVTmark}[1]{\textbf{#1}}
\newcommand{\PEDVTuncertain}[1]{[خوانش نامطمئن: #1]}
\ifnum\PEDVTStandalone=1 \begin{document}\fi

\PEDVTpage{۱}
\begin{itemize}
\item \en{$2/3$} موارد \en{VTE} را \en{DVT} تشکیل می‌دهد و \en{$1/3$} را \en{PE $\pm$ DVT}.
\item در زن و مرد شیوع یکسان دارد.
\item \en{$1/2$} موارد \en{VTE} در بیماران بدون سابقهٔ تروما، جراحی، بی‌حرکتی و بدخیمی رخ می‌دهد.
\item علت عموم مرگ بیماران \en{PE} در بیمارستان:
  \begin{itemize}
  \item \PEDVTmark{\en{PEA}}؛ \PEDVTmark{\en{MOF}} در زمینهٔ نارسایی قلب راست.
  \end{itemize}
\end{itemize}
\subsection*{ریسک فاکتورهای \en{VTE}}
\begin{itemize}
\item سن بالا، بدخیمی، سابقهٔ قبلی \en{VTE}، نارسایی وریدی، بارداری، تروما، بی‌حرکتی، \en{Frailty} (ضعف عمومی).
\item ریسک \en{VTE} در افراد \PEDVTmark{چاق} ۲٫۳، در \en{HTN}، ۱٫۵، در دیابت ۱٫۴، در مصرف سیگار ۱٫۲، در هایپرکلسترولمی ۱٫۲.
\item \PEDVTmark{\en{HF}} در هر دو جنس ریسک \en{VTE} را ۳ برابر می‌کند.
\item ریسک مرگ در \PEDVTuncertain{حاشیهٔ} سه‌ماههٔ بارداری با \en{PE} در سزارین بالاتر است.
\end{itemize}

\PEDVTpage{۲}
\begin{itemize}
\item بارداری، \en{OCP} و هورمون‌درمانی پس از یائسگی ریسک \en{VTE} را \en{$\uparrow$}.
  \begin{itemize}
  \item قرص‌های \en{OCP} که \PEDVTmark{صرفاً پروژسترون} دارند بی‌خطر هستند.
  \end{itemize}
\item زمانی که فردی به‌طور ناگهانی و بدون علت دچار \en{VTE} می‌شود (\en{Unprovoked})، ۵\% احتمال وجود دارد که در او سرطان پیدا شود.
\item درمان \en{Acute VTE} در بیمار سرطانی: \PEDVTmark{\en{NOAC}}.
  \begin{itemize}
  \item دیگر مونوتراپی \en{LMWH} توصیه نمی‌شود.
  \end{itemize}
\item \en{VTE} در زمینهٔ سرطان‌ها شایع است.
\item \en{VTE} در کنسرهای \PEDVTmark{\en{Solid Tumor}}، از جمله آدنوکارسینوم‌ها شایع‌تر است؛ از جمله پانکراس، معده، ریه، مری، پروستات و کولون.
\end{itemize}
\subsection*{\en{Hypercoagulable State}}
\begin{itemize}
\item دو بیماری ژنتیکی با ترومبوفیلی همراه هستند:
  \begin{itemize}
  \item \PEDVTmark{فاکتور \en{V} لیدن}.
  \item \PEDVTmark{موتاسیون ژن پروترومبین}.
  \end{itemize}
\item فاکتور \en{V} لیدن: خطر ۳ برابری \en{VTE}، سقط‌های مکرر.
  \begin{itemize}
  \item اگر \en{OCP} مصرف کند، خطر آن ۱۰ برابر می‌شود.
  \end{itemize}
\end{itemize}

\PEDVTpage{۳}
\begin{itemize}
\item \en{APS}: شایع‌ترین ترومبوفیلی اکتسابی است.
  \begin{itemize}
  \item ترومبوز شریانی و وریدی.
  \item همراهی با سقط‌های مکرر.
  \item اختلالات نورولوژیک حاد.
  \item تشخیص آزمایشگاهی: یکی از \en{Ab}ها حداقل به مدت \PEDVTmark{۱۲ هفته} مثبت باشد.
    \begin{itemize}
    \item آنتی‌کاردیولیپین \en{IgM} یا \en{IgG}.
    \item آنتی بتا ـ ۲ گلیکوپروتئین \en{I}.
    \item آنتی‌پروترومبین.
    \item لوپوس آنتی‌کوآگولان.
    \end{itemize}
  \end{itemize}
\item \en{OCP} و بارداری سطح نرمال پروتئین \en{S} را پایین می‌آورد.
\item \PEDVTmark{\en{CKD}} نیز با افزایش خطر \en{VTE} همراه است.
\end{itemize}
\subsection*{طبقه‌بندی ریسک بیماران \en{PE}}
\begin{itemize}
\item \en{High Risk PE} (شدید):
  \begin{itemize}
  \item ۵–۱۰\% تمام موارد \en{PE}.
  \item مستعد شوک کاردیوژنیک و \en{MOF}.
  \item نارسایی کلیه، ایسکمی کبد، \en{MSC} یافته‌ها شایع هستند.
  \item \en{$1/3$} موارد می‌میرند.
  \item درمان پیشرفته + آنتی‌کوآگولان.
  \end{itemize}
\end{itemize}

\PEDVTpage{۴}
\begin{itemize}
\item \PEDVTmark{دیس‌پنه شایع‌ترین و اصلی‌ترین یافته است}.
\item \en{CP}، سیانوز گذرا.
\item افت فشار نیازمند به وازوپرسور شایع است.
  \begin{itemize}
  \item دادن مایعات زیاد نارسایی قلب راست را تشدید می‌کند.
  \end{itemize}
\item \en{Intermediate Risk PE}:
  \begin{itemize}
  \item \PEDVTmark{فشار خون نرمال است}.
  \item خودش به دو دسته تقسیم می‌شود: \en{Intermediate high}، \en{Intermediate low}.
  \item \en{high}: درمان پیشرفته + آنتی‌کوآگولان.
  \item \PEDVTmark{هم هایپوکینزی قلب راست و هم بیومارکر مثبت}.
  \item \en{Trop}، \en{BNP}، \en{NT-proBNP}.
  \item \en{low}: یکی از موارد بالا وجود داشته، نه هر دو باهم.
  \item در اکو \en{RVE} و افزایش \en{PAP} شایع است.
  \item پایش موارد ضروری است؛ ریسک بالا برای \en{PE} شدید دارند.
  \item آنتی‌کوآگولان.
  \end{itemize}
\item \PEDVTmark{\en{Low Risk PE}}: \PEDVTmark{آنتی‌کوآگولان}.
  \begin{itemize}
  \item \en{BP} طبیعی، عملکرد قلب راست \en{NL} و بیومارکرها منفی.
  \end{itemize}
\end{itemize}

\PEDVTpage{۵}
\subsection*{انفارکت ریوی}
\begin{itemize}
\item \en{CP} پلورتیک، بعضاً هموپتزی.
\item انفارکت ۳–۷ روز پس از \en{PE} رخ می‌دهد.
\item یافته‌ها: \PEDVTmark{تب، لکوسیتوز، \en{$ESR\uparrow$}}.
\item در \en{CXR} شواهد \en{Wedge-shape} و پلورال افیوژن دیده می‌شود.
\end{itemize}
\subsection*{آمبولی پارادوکسیکال}
\begin{itemize}
\item استروک در بیمار با \en{patent foramen Oval} یا \en{Atrial Septal defect}.
\item اغلب در سونو اندام تحتانی \en{DVT} پیدا نمی‌شود.
\end{itemize}
\subsection*{\en{Nonthrombotic PE}}
\begin{itemize}
\item به‌طور کلی ناشایع است.
\item چربی، تومور، هوا، مایع آمنیوتیک.
\end{itemize}
\subsection*{\en{DVT}}
\begin{itemize}
\item \en{Lower Extremity DVT}.
\item \en{Upper Extremity DVT}.
\end{itemize}
\subsection*{\en{Lower Ext DVT}}
\begin{itemize}
\item از اندام فوقانی ۱۰ برابر شایع‌تر است.
\item هرچه پروگزیمال‌تر باشد، احتمال \en{PE} بالاتر است.
\end{itemize}

\PEDVTpage{۶}
\subsection*{\en{DVT} ایزوله در \en{Calf} (ساق)}
\begin{itemize}
\item درمان آنتی‌کوآگولان در موارد علامت‌دار باید انجام شود.
\end{itemize}
\subsection*{\en{DVT} اندام فوقانی}
\begin{itemize}
\item تا حد ممکن کاتتر با تعداد لومن کمتر و قطر کمتر استفاده شود.
\end{itemize}
\subsection*{تشخیص \en{PE}}
\begin{itemize}
\item شایع‌ترین سیمپتوم دیس‌پنه و شایع‌ترین ساین تاکی‌کاردی است.
\item سنکوپ، سیانوز: پروگنوز بد.
\item شواهد به نفع \en{PE} خطرناک:
  \begin{itemize}
  \item کور پولمونال (نارسایی قلب راست) حاد.
  \item \en{JVP} برجسته، گالوپ \en{$S_3$} (در سمت راست)، تاکی‌کاردی، تاکی‌پنه، \en{heave} بطن راست.
  \item شواهد \en{RV dilation} و هایپوکینزی در اکو، یا \en{EKG} با شواهد کور پولمونال: \en{$S_1Q_3T_3$}، \en{New RBBB}، \en{TWI: inf}، \en{TWI $V_1$--$V_4$}.
  \end{itemize}
\item سیمپتوم‌ها: دیس‌پنه، \en{CP} پلورتیک، \en{CP} رتروسترنال، اضطراب، سرفه، هموپتزی، درد و تورم پا.
\item ساین‌ها: تاکی‌پنه، تاکی‌کاردی، تب خفیف، \en{JVP} برجسته، مرمر \en{TR}، اریتم، ادم و تندرنس پا.
\end{itemize}

\PEDVTpage{۷}
\subsection*{سطح \en{D-dimer} بالا}
\begin{itemize}
\item حساس هست ولی اختصاصی نیست.
\item تا یک هفته پس از هر جراحی بالا است.
\item در \en{MI}، سپسیس، کنسر و هر بیماری سیستمیک بالا است.
\item کاربرد صرفاً بیماران سرپایی:
  \begin{itemize}
  \item بیماران اوتپیشن با شک \en{PE} که بیماری سیستمیک دیگری ندارند.
  \end{itemize}
\item زمانی که شک بالینی بیمار قوی است یا اینکه بیمار بستری می‌باشد نباید از \en{D-dimer} استفاده کرد: \en{CT} مستقیم.
\item حداکثر نرمال \en{500 ng/mL}؛ بالای ۵۰ سال: \en{5000 ng/mL}.
% The handwritten value is 5000 ng/mL; it has not been medically corrected.
\item کرایتریا ولز برای \en{PE}: احتمال بالا \en{$>4$}، غیر احتمال بالا \en{$\leq4$}.
% The word above the overwritten label reads احتمال.
\end{itemize}
\subsection*{\en{EKG}}
\begin{itemize}
\item کمک می‌کند \en{MI} و پری‌کاردیت \en{R/o} شود.
\item \en{$S_1Q_3T_3$}، \en{TWI $V_1$--$V_4$}، سینوس تاکی‌کاردی.
\end{itemize}
\subsection*{\en{CXR}}
\begin{itemize}
\item \PEDVTmark{یک عکس \en{near NL} \PEDVTuncertain{در شک بیمار احتمال تشخیص‌شده} قویاً به نفع \en{PE} می‌باشد}.
\item آمبولی پروکسیمال: \en{Focal Oligemia (Westermark Sign)}.
\item \en{Peripheral Wedge shaped density (Hampton hump)}: انفارکت ریه.
\end{itemize}

\PEDVTpage{۸}
\begin{itemize}
\item \en{POCUS}: \en{Consolidation} ساب‌پلورال پارانشیمال؛ در ۷۰\% بیماران.
\item به دلیل وجود آرتیفکت‌ها \en{CT} ممکن است موجب \en{Over-dx} شود.
\item در نمای \en{4-chamber} در \en{CT}:
  \begin{itemize}
  \item بررسی \en{RVD} با نسبت \en{$RV/LV$}: \en{end diastolic diameter Ratio}؛ \en{$>0.9$} دایلاتیشن بطن.
  \item \en{RV/LV Volume Ratio}.
  \item \en{Reflux of the contrast into the IVC}.
  \item \en{InterVentricular Septal bowing toward the LV}.
  \end{itemize}
\item قدم اول درمان \en{PE} طبقه‌بندی ریسک است.
\item [خط‌خورده: \en{PERT}].
\item سه جزء اصلی طبقه‌بندی ریسک:
  \begin{itemize}
  \item[\en{a)}] ارزیابی بالینی بیمار: \en{EKG}، \en{CXR}، اکو، \en{CT} و بیومارکر.
  \item[\en{b)}] ارزیابی سایز و عملکرد \en{RV}.
  \item[\en{c)}] ارزیابی بیومارکرها برای رد انفارکت \en{RV}.
  \end{itemize}
\item برای هر بیمار \en{PE} باید \en{PESI} حساب شود.
\item درمان اصلی \en{PE} آنتی‌کوآگولان است.
\end{itemize}

\PEDVTpage{۹}
\begin{itemize}
\item \en{NOACs}: شروع اثر سریع دارند و ظرف چند ساعت اثر ضدانعقادی شروع می‌شود.
\item اگر قرار به جراحی باشد، از ۲ روز قبل باید قطع شود.
\item به اندازهٔ وارفارین مؤثر ولی از وارفارین ایمن‌تر.
\item اگر آنتی‌کوآگولان قطع شود، خطر \en{VTE} مجدد بالا است.
  \begin{itemize}
  \item ظرف ۱۰ سال ۴۰\% مجدد \en{VTE} می‌کنند؛ \en{$F<M$}.
  \end{itemize}
\item اگر ۶ ماه پس از \en{PE}، \en{CT} مجدد بگیریم، در ۵۰\% بیماران ترومبوز باقی‌مانده است.
\item برای اینکه بدانیم بیمار تا آخر عمر دارو بدهیم، دیگر از پرووکد و آنپرووکد استفاده نمی‌کنیم.
  \begin{itemize}
  \item اگر ریسک \en{PE} مجدد زیر ۳\% در سال باشد می‌توانیم دارو را قطع کنیم.
  \end{itemize}
\item آنچه مهم: \en{Unprovoked, Persistent RF, $\uparrow$D-dimer after 3--6 month}.
\item برای درمان تا آخر عمر فقط \PEDVTmark{\en{NOAC}}: وارفارین و آسپرین نمی‌دهیم.
\end{itemize}
\subsection*{\en{Advanced Therapy} برای \en{PE}}
\begin{itemize}
\item برای بیماران \en{high Risk}.
\item \en{Intermediate high Risk}.
\end{itemize}

\PEDVTpage{۱۰}
\begin{itemize}
\item نشانهٔ موفقیت درمان:
  \begin{itemize}
  \item کاهش \en{RV pressure Overload}.
  \item عدم آزاد شدن سروتونین.
  \end{itemize}
\item دوز ترومبولیتیک:
  \begin{itemize}
  \item \en{Full}: تأیید \en{FDA}.
  \item \en{Half}: \en{off-label}.
  \end{itemize}
\item آلتپلاز: \en{100 mg} وریدی طی ۲ ساعت.
  \begin{itemize}
  \item تا ۱۴ روز از شروع علائم.
  \end{itemize}
\end{itemize}
\subsection*{دوز آنتی‌کوآگولان پروفیلاکسی}
\begin{itemize}
\item بیمار داخلی و جراحی:
  \begin{itemize}
  \item \en{UFH}: \en{5000 u TDS} یا \en{BD}.
  \item \en{Enoxa}: \en{40 mg daily}.
  \end{itemize}
\item بیمار جراحی ماژور ارتوپدی:
  \begin{itemize}
  \item \en{Enoxa}: \en{30 BD} یا \en{40 qd}.
  \item ریواروکسابان: \en{10 qd}.
  \item آسپرین: \en{81 BD}.
  \item ۵ روز اول ریوا، سپس آسپرین روزانه.
  \item آپیکسابان: \en{2.5 BD}.
  \end{itemize}
\end{itemize}
\subsection*{معیار \PEDVTmark{\en{Padua}}}
\begin{itemize}
\item آیا پروفیلاکسی \en{DVT} بدهیم یا نه.
\item \PEDVTmark{امتیاز \en{$\geq4$}}: باید بدهیم.
\end{itemize}

% Repeated subject/date labels and all handwritten page numbers are reproduced.
% The printed notebook icon contains no text.
\ifnum\PEDVTStandalone=1
  \let\PEDVTEnd\enddocument
\else
  \let\PEDVTEnd\relax
\fi
\PEDVTEnd
